\section{Countability and the least counterexample}
\label{sec:countability}

\begin{proposition}[B1: the checked sphere criterion and its extension]
\label{B:spherecriterion}
Let $X$ be a connected CW complex. Suppose that for an infinite set
$D\subseteq\{2,3,\ldots\}$ there are pointed maps
$s_n:S^n\to X$, $f_n:X\to S^n$, and elements
$\alpha_n\in\pi_n(S^3)$ such that
$\deg(f_ns_n)\alpha_n\ne0$. Then
$|[X,S^3]|\geq2^{\aleph_0}$. If $H^1(X;\Z)=0$, then also
$|[X,S^2]|\geq2^{\aleph_0}$.
In particular a return degree prime to the finite order of a nonzero
$\alpha_n$ suffices.
\end{proposition}
\begin{proof}
Choose the base point as a zero-cell and use pointed cellular approximation
for $f_n$ and $s_n$. Give $S^n$ its CW structure with one zero-cell and
one $n$-cell. Thus $f_n$ is constant on $X^{(n-1)}$ and
$s_n(S^n)\subseteq X^{(n)}$. Represent $\alpha_n$ by a pointed map
$a_n:S^n\to S^3=\mathrm{SU}(2)$, and set $g_n=a_nf_n$.
For $A\subseteq D$ set $g_A(x)=\prod_{n\in A}g_n(x)$ in increasing
order. On each closed cell only finitely many factors are nonconstant.
The weak topology of a CW complex therefore proves continuity.

If $A\ne A'$, take the least $n$ in their symmetric difference. On
$X^{(n)}$ the two products have the form $h g_n$ and $h$, in one
order or the other. A homotopy between $g_A$ and $g_{A'}$ would,
after restriction and pointwise left multiplication by $h^{-1}$, nullhomotope
$g_n|_{X^{(n)}}$. Composing with $s_n$ contradicts
$[a_nf_ns_n]=\deg(f_ns_n)\alpha_n\ne0$. Pointed and free classes
agree for the simply connected target $S^3$. This injects the power set
of a countably infinite subset of $D$ into $[X,S^3]$.

For the second assertion let $p:S^3\to S^2$ be the Hopf principal
$S^1$-bundle. If $p u$ and $p v$ are homotopic, lift that homotopy
starting at $u$. Its endpoint and $v$ are lifts of the same map and
differ by a continuous function $X\to S^1$. This function is
nullhomotopic because $[X,S^1]=H^1(X;\Z)=0$. Hence $u$ and $v$
are homotopic. Thus $p_*:[X,S^3]\to[X,S^2]$ is injective.
This checks every step of the supplied Proposition~4.1 and proves
the requested weaker hypothesis.
\end{proof}

\begin{proposition}[B1a: finite domination]
\label{B:finitedomination}
If $X$ is homotopy dominated by a finite CW complex and $Y$ has
countably many components and countable homotopy groups, then
$[X,Y]$ is countable. In particular the finite-domination remark in R4
is correct for $Y=S^2$.
\end{proposition}
\begin{proof}
Let $i:X\to K$ and $r:K\to X$ satisfy $ri\simeq\id_X$, with
$K$ finite. Precomposition with $r$ injects $[X,Y]$ into $[K,Y]$,
since $i^*r^*=\id$. The latter set is countable: choose the image
components, then extend maps cell by cell. At each of the finitely many
steps the choices, when they exist, form a quotient of a countable
homotopy group (or a finite product of such groups). The fundamental
group choices on the finite one-skeleton are also countable.
\end{proof}

\begin{theorem}[B2--B3: the least counterexample]
\label{B:main}
Let $M$ be the six-element monoid specified in Part~A. Then
\[
 |[BM,S^3]|=|[BM,S^2]|=2^{\aleph_0}.
\]
Consequently the functor $C\mapsto[BC,S^2]$ on finite categories
cannot be represented as a sequential colimit of representables by
finite categories. Six is the least possible number of arrows in a
counterexample.
\end{theorem}
\begin{proof}[Proof of countability assertion; minimality is proved below]
The explicit wedge equivalence established in Part~A gives, for any
simply connected target $Y$,
\[
 [BM,Y]\cong\prod_{n\geq2}\pi_n(Y).
\]
The bijection is restriction to the sphere summands; its inverse glues
pointed maps, and homotopies glue by the CW weak topology.
For $Y=S^3$ infinitely many factors are nonzero by Serre's theorem
\cite[Theorem~10 and the sphere corollary, pp.~217--219]{CDSerre1953}.
All factors are countable, so the product has exactly
the cardinality of the continuum. The Hopf fibration identifies
$\pi_n(S^3)$ with $\pi_n(S^2)$ for $n\geq3$, and adds the
countable factor $\pi_2(S^2)=\Z$, giving the same cardinality.
Equivalently one may apply Proposition~\ref{B:spherecriterion} and the countable upper bound
on each homotopy group. A sequential colimit of the finite sets
$\pi_0\Fun(C,K_i)$ is countable for any fixed finite $C$.
Taking $C=M$ proves the obstruction directly, without needing the
converse direction of R1.
\end{proof}

\begin{lemma}[B3a: monoids of order at most five]
\label{B:smallmonoids}
For every monoid $L$ with at most five elements, $BL$ is homotopy
equivalent either to the classifying space of a finite group or to $S^2$.
\end{lemma}
\begin{proof}
The minimal nonempty ideal of a finite monoid is a completely simple
semigroup, hence a Rees matrix semigroup $I\times H\times J$.
If $|I|=1$ or $|J|=1$, its group $H$ is a homotopy equivalent
corner: there is an idempotent $e$ with $eLe=eL$ or $eLe=Le$.
In the former case $l\mapsto el$ is a unital homomorphism $L\to H$
(the identity of $H$ is $e$). In the Karoubi category on $1,e$,
the component $e:1\to e$ gives a natural transformation from the
object-$1$ inclusion to this homomorphism followed by the object-$e$
inclusion: its naturality equation is $(el)e=el$. Restriction to
$H$ is the identity, and the object-$1$ inclusion is a classifying-space
equivalence by the comma-category proof in Part~A. The other case
is its opposite. This is
the direct proof of Sweeney's Theorem~2.7 and Corollary~2.8.

If both $|I|,|J|\geq2$, the minimal ideal has at least four elements.
The identity is outside that ideal: otherwise the monoid is simple,
and the identity forces both index sets to be singletons. The order
bound therefore forces $|I|=|J|=2$, $H=1$, and $L=1\sqcup(I\times J)$.
This is the rectangular-band monoid $P_2$, whose classifying space
is $S^2$ by the explicit Part~A proof with no $y$-letters.
\end{proof}

\begin{lemma}[B3b: categories of at most five arrows]
\label{B:smallcategories}
If $C$ has at most five arrows, $[BC,S^2]$ is countable.
\end{lemma}
\begin{proof}
It is enough to consider connected $C$, since there are finitely
many components. One object is covered by Lemma~\ref{B:smallmonoids} and Miller's
Theorem~A, which says that $\operatorname{Map}_*(BH,S^2)$ is weakly
contractible for a finite group $H$
\cite[Theorem~A, p.~39]{Miller1984,Miller1985Correction}.

With two objects and arrows in both directions, choose $f:x\to y$
and $g:y\to x$. If both composites were nonidentity, the two
identities, $f,g,gf,fg$ would already give six arrows. Hence one
object is a retract of the other. The inclusion of the latter's full
endomorphism category is a homotopy equivalence: it becomes an
equivalence after splitting idempotents, and idempotent completion
preserves the classifying space. Its endomorphism monoid has at most
three elements and is covered by Lemma~\ref{B:smallmonoids}.

Otherwise all cross-arrows have one direction; let their number be $c$.
If $c=1$, the category is the categorical join of its two endomorphism
monoids. If either has just an identity its classifying space is a cone.
The only other possibility within the bound is two order-two monoids.
Each is either the group $C_2$ or the idempotent monoid, whose
classifying space is contractible. The only remaining case is
$BC_2*BC_2$. This is the homotopy pushout of the two projections
$BC_2\leftarrow B(C_2\times C_2)\to BC_2$. Applying derived mapping
spaces into $S^2$, Miller's theorem makes all three unpointed mapping
spaces weakly equivalent to $S^2$ via the constant maps. The induced
homotopy pullback has the same components as $S^2$, so there is only
one free homotopy class.

If $c=2$, there is at most one extra endomorphism. Without it the
nerve is a circle. With it, the nontrivial endomorphism acts on the
two cross-arrows, either as an involution (group $C_2$) or as an
idempotent transformation. The resulting nerve is the homotopy
cofiber of
\[
 B(S\mathbin{/\!/}L)\longrightarrow BL,
\]
where $S$ is the two-element cross-arrow set and
$S\mathbin{/\!/}L$ is its action category. This follows directly
by partitioning nerve simplices into those at the two objects and
those with their unique cross-arrow; the mixed simplices are the
two-sided bar model for that mapping cylinder.
For the group, a transposition gives a contractible action category,
and trivial action gives two copies of $BC_2$. The resulting types
are $BC_2$ and $S^1\vee\Sigma BC_2$, respectively. For the
idempotent monoid, $BL$ is contractible; constant action has
contractible action category and identity action has two contractible
components. The resulting types are a point and $S^1$.
Miller's theorem also gives $[\Sigma BC_2,S^2]_*=0$ through the
first homotopy group of $\operatorname{Map}_*(BC_2,S^2)$.
Thus all these cases have one class. If $c=3$, there are no extra
endomorphisms and the type is $S^1\vee S^1$, again with one class.

Finally, a connected category on $k\geq3$ objects needs at least
$k-1$ nonidentity arrows, so $2k-1\leq5$ forces $k=3$ and two
nonidentity arrows. They form a tree. They cannot be composable,
since their nonidentity composite would be a third arrow. Thus their
common source or common target is initial or terminal, and the nerve
is contractible.

The independent exhaustive program \texttt{B\_small\_categories.cpp}
checks all composition tables within this bound, with nonidentity
arrow labels normalized into source--target blocks, and verifies exactly
these structural alternatives. Together with the six-arrow example,
this proves the claimed optimality.
\end{proof}
